\ifdefined\pdftrailerid\pdftrailerid{}\fi
\documentclass[11pt]{article}
\usepackage{amsmath,amssymb,amsthm,hyperref}
\usepackage[margin=1in]{geometry}
\newtheorem{theorem}{Theorem}
\newtheorem{lemma}[theorem]{Lemma}
\DeclareMathOperator{\des}{des}
\title{Real zeros of descent polynomials for permutations with a minimum cycle length}
\author{Bowen Zuo\\ {\small East China University of Science and Technology}}
\hypersetup{pdfauthor={Bowen Zuo},pdftitle={Real zeros of descent polynomials for permutations with a minimum cycle length}}
\date{26 November 2026}
\begin{document}
\maketitle
\begin{abstract}
Let $D_{n,k}$ enumerate permutations of $n$ letters with all cycles of length at least $k$, by ordinary descents. We give a proof of the real-rootedness conjecture for $3\le k\le n$. Except at $n=5$, a single system of disjoint negative intervals at each size locates one simple root of every $D_{n,k}/t$ per interval. At large sizes we obtain the intervals by bounding the lower-degree terms of the Gessel--Reutenauer alphabet polynomial, then using two different estimates for alternating power sums. The finite range is covered by reproducible integer sign certificates. We also prove real-rootedness of every nonzero nonnegative combination at fixed size, including the ordinary Eulerian and derangement rows, and classify its multiple roots.
\end{abstract}

\section{The statement}
Write $\mathcal A_{n,k}$ for the permutations of $[n]$ whose shortest cycle has length at least $k$. Put
\[
 D_{n,k}(t)=\sum_{\sigma\in\mathcal A_{n,k}}t^{\des(\sigma)},\qquad
 \des(\sigma)=\#\{i<n:\sigma(i)>\sigma(i+1)\}.
\]
Alexandersson asks whether these polynomials have only real roots \cite{open}. The case $k=1$ is Eulerian and Alexandersson proved real-rootedness for the ordinary descent derangement row $k=2$ \cite[Theorem~2.1]{der}. Theorem~\ref{allcuts} concerns the compatibility of all the cutoffs at a fixed size. Until the final section $3\le k\le n$. All such rows are nonempty because they include the $n$-cycles.

ProgramComputer~\cite{pc} proved the single-cycle case at every size,
including the double root at $n=5$. For $n\ge25$ and $k\ge52\log n$
they also proved simple nonpositive zeros for arbitrary nonzero
nonnegative mixtures of the cycle types with minimum part at least $k$.
Our mixtures are of the whole cutoff rows, and the result covers every
$1\le k\le n$.

\begin{theorem}\label{main}
For each $n\ge3$, $D_{n,k}$ has degree $n-2$ and a simple root at zero. For $n\ne5$ there are $n-3$ disjoint open intervals in $(-\infty,0)$, independent of $k$, each containing one simple root of every $D_{n,k}$. For $n=5$ all three rows equal $6t(1+t)^2$.
\end{theorem}
In particular any nonzero nonnegative linear combination at a fixed size has the same root counts and multiplicities. At $n=5$ the combination is a positive multiple of $6t(1+t)^2$. For $n\ne5$ the endpoint signs are strict and common to the whole family, and the zero at the origin supplies the remaining root. At $n=3$ the interval collection is empty. Enumeration by the descent set of a single cycle is developed in \cite{et}, and plethystic formulas for cycle-type refinements in \cite{gz}. The other long-cycle conjecture in \cite{open} concerns excedances and interlacing as $n$ changes.

\section{The alphabet polynomial}
Let $M_j(q)=j^{-1}\sum_{d\mid j}\mu(d)q^{j/d}$. Define
\begin{equation}\label{letter}
 f_n(q)=[z^n]\prod_{j\ge k}(1-z^j)^{-M_j(q)},\qquad
 \frac{D_{n,k}(t)}{(1-t)^{n+1}}=\sum_{q\ge1}f_n(q)t^{q-1}.
\end{equation}
The second identity is the finite-alphabet specialization of Gessel and Reutenauer \cite[Theorem~3.6]{gr}, with the first $q$ variables set to one and the rest to zero. To see its normalization, a permutation with $r$ descents admits $\binom{q+n-1-r}{n}$ weakly increasing $q$-letter labelings with strict rises at the descents. The cycle itineraries of these labeled permutations give multisets of primitive necklaces. There are $M_j(q)$ necklace types of length $j$, each available with arbitrary multiplicity. Summing the binomial coefficients in $q$ gives $t^r/(1-t)^{n+1}$.

At $q=0$ every $M_j$ vanishes, so the product is one and $f_n(0)=0$ for $n>0$. Write $f_n(q)=\sum_{h=1}^n a_h q^h$ and $p_n=a_n$. Only the leading terms with cycle multiplicity one in the logarithm of the product contribute at degree $n$. Thus
\[
 p_m=[z^m]\exp\Big(\sum_{j\ge k}z^j/j\Big),\qquad
 m p_m=\sum_{s=0}^{m-k}p_s\quad(m\ge k).
\]
Here $p_m$ is the fraction of admissible permutations, $p_0=1$ and $p_s=0$ for $0<s<k$. For $m\ge k$ we have $1/(2k)\le p_m\le1/k$. Both bounds hold for $k\le m<2k$, where $p_m=1/m$. Induction for $m\ge2k$ gives a numerator between $1+(m-2k+1)/(2k)$ and $1+(m-2k+1)/k$.

The reversed row $t^{n-1}D_{n,k}(1/t)$ has alphabet polynomial $f_n^*(q)=(-1)^nf_n(-q)$. This follows from
$\binom{-q+n-1-r}{n}=(-1)^n\binom{q+r}{n}$. Both alphabet polynomials vanish at $q=1$, since the increasing and the decreasing permutations are excluded. In particular $D_{n,k}(0)=0$ and $\deg D_{n,k}\le n-2$.

\section{Uniform control of the lower degrees}
Put $E(v)=p_n^{-1}\sum_{h<n}|a_h|v^{n-h}$ and $V=2/3$. Substituting $zv$ for $z$ and $1/v$ for $q$ in the logarithm of \eqref{letter} splits off the terms $\sum_{j\ge k}z^j/j$. The remaining terms are
\[
 B=\sum_{j\ge k}\sum_{b\ge1}
 \sum_{\substack{d\mid j\\(b,d)\ne(1,1)}}
 \frac{\mu(d)}{jb}z^{jb}v^{jb-j/d}.
\]
At $0<v\le V$ the total absolute mass at $z=1$ is at most
\[
 J(v)=v^2/3+v^2/4+v^4/5+(v^3+v^4+v^5)/6+v^6/7
       +\frac{v^4}{1-v}+\frac{v^3}{2(1-v)(1-v^3)}.
\]
For $b=1$ the first five terms account for $j=3,\ldots,7$. When $j\ge8$ the number of divisors other than one is at most $j/2$ and the exponent is at least $\lceil j/2\rceil$. For $b\ge2$ bound $1/b$ by $1/2$, use at most $j$ divisors and sum $v^{j(b-1)}$ over $j\ge3$.

Let $C(z,v)$ replace the coefficients of the displayed bad logarithm $B$ by their absolute values and put $H=e^C-1$. The convolution with the good part gives $E(v)\le p_n^{-1}\sum_{r=1}^n p_{n-r}[z^r]H(z,v)$. A term with remaining length $s>0$ is multiplied by $p_s/p_n\le2$. Terms with $s=0$ come entirely from the error exponential. Each of their factors loses at least half its $z$-degree, since $d\ge2$ or $b\ge2$. Their loss is at least $n/2$, and $1/p_n\le2n$. Consequently
\[
 E(v)\le 2(e^{J(v)}-1)+2n(v/V)^{n/2}e^{J(V)},\qquad E(V)\le2n e^{J(V)}.
\]
For the second bound use $p_s\le1$ directly for every remaining length. For $n\ge128$, $J(V)<2$, $e^2<8$, and $16n(3/4)^{n/2}<1/1000$. At the four smaller rational arguments, use $e^J-1\le J/(1-J)$ to obtain
\begin{equation}\label{est}
 E(1/4)<1/7,\quad E(1/3)<1/3,\quad E(3/8)<3/7,\quad E(1/2)<2 .
\end{equation}
The tightest margin is \[ \frac37-\frac{2J(3/8)}{1-J(3/8)}-\frac1{1000} = \frac{18913056563}{25670690437000}>0. \] The exponential bound in $n$ decreases from its value at 128 by the ratio test.

\section{A common sign at each node}
Assume $n\ge128$ throughout this section. Set $N=n+1$, $\theta=\pi i/N$, $x=\pi\cot\theta$ and $t_i=-e^{-x}$ for $2\le i\le n-1$. It suffices to show that the sign of $D_{n,k}(t_i)$ is $(-1)^{i-1}$. We first take $i\le N/2$. The evaluation has the sign of
$\sum_{q\ge1}(-1)^{q-1}f_n(q)e^{-xq}$, with the Abel limit at $x=0$. The same estimates below apply with $f_n^*$ in place of $f_n$.

\paragraph{Central nodes.}
When $i\ge\sqrt N/2$, write $s=\sin\theta$ and $\rho=\pi/s$. For $h\ge1$ the alternating $h$th power sum equals
\[
 -h!\sum_{\ell\text{ odd},\,\ell\in\mathbb Z}(x+\mathrm i\pi\ell)^{-h-1}.
\]
Here $\mathrm i^2=-1$. The logarithmic derivative of the product for $\cosh(x/2)$ gives
\[
 (e^x+1)^{-1}=\tfrac12-\sum_{\ell\ge1\ {\rm odd}}2x/(x^2+\pi^2\ell^2).
\]
Differentiate $h$ times and multiply by $(-1)^h$. The paired expansions and their derivatives converge uniformly on compact real intervals, which includes the Abel endpoint.
For $h=n$ the nearest pair contributes $(-1)^{i-1}A$, where $A=2n!\rho^{-N}$. We have $Ns^2\ge2$, and the remaining odd pairs contribute in absolute value at most $A/4$. Indeed Bernoulli's inequality bounds the normalized sum by $\sum_{\ell\ge3\ {\rm odd}}(\ell^2-1)^{-1}=1/4$.

For $n\ge128$, $\sin(\pi/(2\sqrt N))\ge(199/200)\pi/(2\sqrt N)$ by $\sin z\ge z-z^3/6$. Thus $\rho\le(41/20)\sqrt n$. If $n/2\le h<n$, the sum of relative moduli of the positive odd poles at exponent $h+1$ is at most $3/2$, since $(h+1)s^2/2\ge1/2$. Also $h!/n!\le(2/n)^{n-h}$ and $2\rho/n\le3/8$. By \eqref{est} these terms contribute at most $(9/14)p_nA$. For $h<n/2$, the geometric mean of the last $n-h$ factorial factors is at least $n/e$. The sum of relative pole moduli is at most $\sum_{\ell\ge1\ {\rm odd}}\rho^2/(\pi\ell)^2=\rho^2/8\le n$. Since $e\rho/n\le1/2$, the contribution divided by $p_n A$ is at most
$n(3/4)^{n/2}E(V)\le16n^2(3/4)^{n/2}<1/100$. The last quantity decreases for $n\ge128$. The total relative error is below $1/4+9/14+1/100<1$.

\paragraph{End nodes.}
Let $i<\sqrt N/2$ and $u_q=q^n e^{-xq}$. Its successive ratios decrease in $q$. From
$N/i-\pi^2 i/(2N)\le x\le N/i$ (use $\cos\theta\ge1-\theta^2/2$, $\sin\theta\le\theta$ and $\tan\theta\ge\theta$), one obtains
\[
 \log(u_{i+1}/u_i)\le 1/(4i)-n/(2i(i+1)),\qquad
 \log(u_{i-1}/u_i)\le 1/i-n/(2i^2).
\]
Here $\pi^2 i/(2N)<5/(4i)$, $\log(1+1/i)\le1/i-1/(2i(i+1))$ and $\log(1-1/i)\le-1/i-1/(2i^2)$. Both ratios are less than $1/4$. For $i\ge4$ substitute $n\ge4i^2-1$ in the two bounds; their maxima are at most $-121/80$ and $1/4-(4-1/16)/2$. For $i=2,3$ use $n\ge128$.

Write $g=f_n$ or $f_n^*$. For $q\ge4$, \eqref{est} places $g(q)/(p_n q^n)$ between $6/7$ and $8/7$. At $q=3$ the bounds are $2/3$ and $4/3$, and at $q=2$ the upper bound is 3. For $q=1$ the bound $g(1)/(p_n1^n)\le3$ also holds here, since $g(1)=0$.
For all $i\ge4$, the left ratio $\eta=u_{i-1}/u_i$ satisfies
\[
 \log\eta\le-2+\frac1i+\frac1{2i^2}\le-\frac{55}{32}<-\frac53,\qquad
 e^{5/3}>\sum_{a=0}^4\frac{(5/3)^a}{a!}=\frac{10009}{1944}>5.
\]
Thus the left and right outward ratios are bounded by $1/5$ and $1/4$.
In units of $p_nu_i$, first charge all off-center terms at amplitude $8/7$.
For $q=3$ the excess is at most $4/21$; for $q=2,1$ each excess is at most $13/7$.
The largest resulting bound occurs at $i=4$ and is
\begin{equation}\label{tails}
 \frac87\left(\frac14+\frac13\right)
 +\frac4{21}\frac15+\frac{13}{7}\left(\frac1{25}+\frac1{125}\right)
 =\frac{2084}{2625}<\frac67 .
\end{equation}
The difference $166/2625$ is a strict margin against the center term.
For $i=3$ both neighboring ratios are at most $\exp(1/3-n/24)<1/100$,
so both tails including $q=1$ cost at most $6/99<2/3$.

At $i=2$ use the unrestricted bound $0\le g(q)\le q^n$, obtained by summing labelings over all permutations, also after reversal. Single cycles alone give $g(2)\ge2^n/(2n)$: their contribution is $M_n(2)$ or $(-1)^nM_n(-2)$, and the sum of the absolute lower powers in the divisor formula is at most $n2^{n/2}<2^{n-1}$. Both outward ratios are at most $\beta=\exp(1/2-n/12)$. Here $8n\beta<1$ at 128 and decreases thereafter, so $2\beta/(1-\beta)<1/(2n)$. This proves the desired sign throughout the end region.

Finally $t_{N-i}=1/t_i$. The reversed polynomial contributes the sign factor $(-1)^{n-1}$ when passing to the other half of the mesh. Hence the $n-2$ nodes strictly alternate and yield $n-3$ distinct negative roots by continuity. Together with zero these exhaust the degree bound.

\section{The finite certificate}
For $3\le n\le127$ the companion data give rational dyadic nodes at each size. For each $k$ the integer coefficients are reconstructed by
\[
 [t^r]D_{n,k}=\sum_{j=0}^r(-1)^j\binom{n+1}{j} f_n(r-j+1).
\]
The two independent reconstructions either insert long-necklace factors or remove the short factors from $(1-qz)^{-1}$. The latter equality for the unrestricted product follows from $\sum_{j\mid b}jM_j(q)=q^b$. For $k>n/2$ only $n$-cycles occur and one row represents all those cutoffs. From a node $a/2^s$ the verifier evaluates the numerator by integer Horner after clearing the positive denominator. Checking order and strict alternation costs 327,451 evaluations on the common meshes for the relevant distinct rows outside $n=5$. The exceptional row factors exactly. The degree argument from the preceding section applies to the dyadic nodes without any rounding assumption. This completes the finite part of Theorem~\ref{main}.

\section{Including the Eulerian and derangement rows}

\begin{theorem}\label{allcuts}
For every $n\ge1$ and every $w_1,\ldots,w_n\ge0$ with $\sum w_k>0$,
the zeros of $P(t)=\sum_{k=1}^n w_kD_{n,k}(t)$ are real and nonpositive.
Every zero is simple except in the following two cases. If $n=4$ and only $w_2$ is positive, $-2$ is a double zero. If $n=5$ and $w_1=w_2=0$, $-1$ is a double zero.
\end{theorem}

For $k=1$ the alphabet polynomial in \eqref{letter} is $q^n$ in both orientations.
For $k=2$ remove the fixed-point factor from the unrestricted product:
\[
 f_n(q)=[z^n]\frac{(1-z)^q}{1-qz}=q^nG_n(1/q),\qquad
 G_n(v)=\sum_{j=0}^n\frac{(-1)^j}{j!}\prod_{\ell=0}^{j-1}(1-\ell v).
\]
The ambient reversal $t^{n-1}D_{n,2}(1/t)$ has alphabet input $f_n^*(q)=(-1)^nf_n(-q)=q^nG_n(-1/q)$.
Since $\deg G_n\le n-1$, neither orientation introduces a constant power of $q$.
At $q=1$ the Eulerian inputs are both one; the derangement inputs are
\[
 f_n(1)=0,\qquad f_n^*(1)=\mathbf1_{2\mid n},
\]
since $f_n(-1)=[z^n](1-z^2)^{-1}$ for $n\ge1$.
The coefficient of $q^n$ in either orientation is $p_n=G_n(0)$.

For a series $H(v)$ and $r\ge0$ write $\|H\|_r=\sum_{a\ge0}|[v^a]H|r^a$.
For $V=2/3$ the binomial series $G(v)=(1-v)^{1/v}$, with $G(0)=e^{-1}$, gives
\[
 \|G-G_n\|_V\le(4/5)^n\sum_{j\ge0}
 \frac{(5/4)^j}{j!}\prod_{\ell=0}^{j-1}(1+\ell V)
 = (4/5)^n6^{3/2}<10^{-10}=:\epsilon \quad(n\ge128).
\]
The omitted sum starts at $j=n+1$; the majorant is the binomial series $(1-5V/4)^{-1/V}$. The estimate controls the norm at every smaller radius and evaluation at every $|v|\le V$.
Put $S(v)=\sum_{a\ge1}v^a/(a+1)\le v/(2(1-v))$ for $0\le v\le V$.
Since $G(v)=e^{-1}\exp(-S(v))$, the relative positive-degree coefficient mass of $G_n$ at radius $v$ is at most $\bigl(e^{-1}(e^{S(v)}-1)+\epsilon\bigr)/p_n$.
Alternating exponential sums give $0.3678<p_n<0.3680$ and $0.3678<e^{-1}<0.3680$.
For $0\le b\le1$ the Taylor series gives
$e^b\le1+b+b^2/(2(1-b/3))$. At $v=3/8,1/2,2/3$ the mass bounds are respectively below $3/7,2,3$. They give the required central-node bounds in both orientations.

For the end nodes, $\log G(v)=-\int_0^1(1-vu)^{-1}\,du$ is decreasing.
Let $g=f_n$ or $f_n^*$. For $q\ge4$, $G_n(\pm1/q)$ lies between $(3/4)^4-\epsilon$ and $(5/4)^{-4}+\epsilon$. Dividing by $p_n$ puts $g(q)/(p_nq^n)$ strictly between $6/7$ and $8/7$. At $q=3$ the endpoint values $(2/3)^3$ and $(4/3)^{-3}$ similarly give $2/3$ and $4/3$. At $q=2$ the lower bound used at node $i=2$ is
\[
 \frac{g(2)}{2^n}=G_n(\pm1/2)\ge\frac14-\epsilon>\frac1{2n}.
\]
In units of $p_n2^n$ the corresponding ratio is $G_n(\pm1/2)/p_n$.
For both orientations $0\le g(q)\le q^n$: the nonnegative binomial contributions from derangements form a subfamily of the full Eulerian sum, also after ambient reversal. Thus $g(q)/(p_nq^n)<3$ for every integer $q\ge1$. The Eulerian row has $p_n=1$ and this ratio equals one. Equation~\eqref{tails} now applies for all $i\ge4$, including the $q=1$ term. At $i=3$ the bound $6/99$ applies, and at $i=2$ the preceding lower bound, in units of $2^n$, and $8n\exp(1/2-n/12)<1$ suffice. Reflection gives the same signs on the other half of the mesh.

For $3\le n\le127$, except at $4,5$, the two added rows pass 15,740 exact sign tests on the original rational dyadic nodes. For $1\le q\le n$ the integer derangement counts simplify to $f_n(q)=q^{n-q}(q-1)^q$. A second verifier uses the triangular equations $f_n(q)=\sum_{d=0}^{q-1}[t^d]D_{n,2}(t)\binom{q+n-1-d}{n}$, whose last multiplier is one. For $k=1$ the corresponding input is $q^n$.

For $n\ge3$ the two endpoint identities are
\[
 P(0)=w_1,\qquad [t^{n-1}]P=w_1+\mathbf1_{2\mid n}\,w_2.
\]
Only the identity has no descents; the decreasing permutation has a fixed point for odd $n$ and only transpositions for even $n$.
The $n-2$ strict mesh signs, running from $-$ nearest zero to $(-1)^{n-2}$ farthest left, give $n-3$ distinct negative roots between nodes. If $w_1>0$, there is a further root at each outer side and the degree is $n-1$. If $w_1=0$, zero is a simple root: $(2,3,\ldots,n,1)$ in one-line notation is an $n$-cycle with one descent. When $n$ is even and $w_2>0$, the degree is $n-1$ and there is a further root to the left. Otherwise the degree is at most $n-2$. In each case the distinct roots exhaust the degree.

For $n=4,5$ the distinct rows, in order of cutoff $1,2,\ge3$, are
\[
\begin{array}{c|lll}
4&1+11t+11t^2+t^3&t(t+2)^2&3t(1+t)\\
5&1+26t+66t^2+26t^3+t^4&8t+24t^2+12t^3&6t(1+t)^2 .
\end{array}
\]
At $n=4$ all rows are negative at $-1/2$; at $-2$ they are positive except for the derangement row, which vanishes. The preceding degree count applies whenever a weight outside $k=2$ is positive. At $n=5$ all rows are negative at $-1/4$ and $-4$; at $-1$ their values are $16,4,0$ in the displayed groups. If $w_1+w_2>0$ the same count again gives only simple roots. The displayed factorizations give both remaining double roots. For $n=2$ the mixture is $w_1+(w_1+w_2)t$; for $n=1$ it is a positive constant. This proves Theorem~\ref{allcuts}.


\paragraph{Reproducibility.}
The ancillary directory contains the finite sign data, two independent Python implementations for the added rows, and commands to replay the certificates.
The Lean proofs start from the permutation sum defining $D_{n,k}$. The module \texttt{DescentAllSizes} verifies the common isolating intervals of Theorem~\ref{main} for every $n\ge3$ other than $5$, and real splitting also at $n=5$. The module \texttt{DescentAllCutoffSizes} verifies Theorem~\ref{allcuts} for every $n\ge1$, including the two exceptional supports. The accompanying source package includes the dependency locks and replay instructions. AI tools assisted with derivation, verification, and drafting. The author takes responsibility for the manuscript.

\begin{thebibliography}{9}
\bibitem{open} P.~Alexandersson, Real-rooted polynomials: open problems,
\emph{The symmetric functions catalog}, accessed 26 November 2026,
\url{https://www.symmetricfunctions.com/realRootedOpenProblems.htm#permutations-with-long-cycles}.
\bibitem{der} P.~Alexandersson, Real-rooted Eulerian polynomials from permutations, words, and paths,
arXiv:2609.07325v2 (2026).
\bibitem{gr} I.~M.~Gessel and C.~Reutenauer, Counting permutations with given cycle structure and descent set, \emph{J. Combin. Theory Ser. A} \textbf{64} (1993), 189--215.
\bibitem{et} S.~Elizalde and J.~Troyka, Exact and asymptotic enumeration of cyclic permutations according to descent set, \href{https://doi.org/10.1016/j.jcta.2019.02.012}{DOI 10.1016/j.jcta.2019.02.012}.
\bibitem{gz} I.~M.~Gessel and Y.~Zhuang, Plethystic formulas for permutation enumeration, \href{https://arxiv.org/abs/2001.00654}{arXiv:2001.00654}.
\bibitem{pc} ProgramComputer, Real zeros of descent polynomials for permutations with long cycles, research manuscript, 25 November 2026. \href{https://github.com/ProgramComputer/open-math-research/blob/4da48b72ed8f92f96c2faf60e2d63d253b3b507a/topics/long-cycle-descents/paper.tex}{Archived source}.
\end{thebibliography}
\end{document}
